% Section owner: Manas (architecture and training engine)
\section{Architecture and Training}
\label{sec:p2arch}
\subsection{Limits of the Phase 1 CNN}
The Phase 1 CNN applies three blocks of one $3\times3$ convolution, ReLU and max-pooling, flattens the
$128\times3\times3$ output and classifies with two dense layers. Three properties limit it
(Table~\ref{tab:p2arch}). First, \DenseSharePone\% of its \ParamsPone{} parameters sit in dense layers, which do
not share weights across positions. Second, each resolution receives a single non-linear step, although the
confusable classes differ in small local patterns such as collars, plackets and sleeve length. Third, a unit of its
last feature map sees only $\RFPone\times\RFPone$ of the $28\times28$ pixels.

\subsection{FashionNet}
The Phase 2 family consists of a stem (one $3\times3$ convolution with 32 channels, batch normalisation, ReLU),
three stages of widths 32, 64 and 128 at resolutions 28, 14 and 7 separated by $2\times2$ max-pooling, global
average pooling, dropout and one linear layer. A stage contains $b$ blocks; a block computes
\begin{align}
F(x) &= \mathrm{BN}\bigl(\mathrm{conv}(\mathrm{ReLU}(\mathrm{BN}(\mathrm{conv}(x))))\bigr), \nonumber\\
y &= \mathrm{ReLU}\bigl(F(x) + s(x)\bigr),
\label{eq:p2block}
\end{align}
where $s$ is the identity, or a $1\times1$ convolution with batch normalisation when the number of channels
changes~\cite{he2016deep}. Since $\partial y/\partial x$ contains the identity term from $s$, the gradient reaches
early layers without passing through every convolution. Two stacked $3\times3$ convolutions cover a $5\times5$
region with 18 instead of 25 weights per channel pair~\cite{simonyan2015very}. Global average
pooling~\cite{lin2014network} replaces the $1152\rightarrow128$ dense layer; its 128 outputs are the embedding used
by the hybrid models. No pooling follows the last stage, because $7\times7$ cannot be halved without discarding a
row and a column.

Three switches turn the family into an experiment: \emph{residual} (remove $s$), \emph{batch normalisation}
(replace it by the identity and give the convolutions a bias), and the number of blocks per stage. The base model
uses $b=1$ (7 convolution layers, \ParamsBase{} parameters, receptive field \RFBase{} pixels); the deeper variant
uses $b=2$ (13 layers, \ParamsDeeper{} parameters).

% Generated by report/build_phase2_report_assets.py. Do not edit by hand.
\begin{table*}[t]
\centering
\caption{Architecture variants. Dense share is the fraction of parameters in fully connected layers; the receptive field is that of the last feature map. Validation macro-F1 is the mean $\pm$ standard deviation over three seeds.}
\label{tab:p2arch}
\footnotesize
\begin{tabular}{lrrrrr}
\toprule
Model & 3$\times$3 conv layers & Parameters & Dense share & Receptive field & Val macro-F1 \\
\midrule
Phase 1 CNN & 3 & 241,546 & 61.6\% & 22 px & 0.9343 $\pm$ 0.0023 \\
FashionNet, no shortcut & 7 & 297,450 & 0.4\% & 34 px & 0.9510 $\pm$ 0.0024 \\
FashionNet base & 7 & 308,074 & 0.4\% & 34 px & 0.9523 $\pm$ 0.0007 \\
FashionNet deeper (final) & 13 & 696,042 & 0.2\% & 62 px & 0.9557 $\pm$ 0.0010 \\
FashionNet narrow & 7 & 77,754 & 0.8\% & 34 px & 0.9413 $\pm$ 0.0014 \\
FashionNet, no BatchNorm & 7 & 307,402 & 0.4\% & 34 px & 0.9453 $\pm$ 0.0050 \\
\bottomrule
\end{tabular}
\end{table*}


\subsection{Training recipe}
The Phase 1 recipe is Adam with a constant learning rate of $10^{-3}$, no augmentation and early stopping with
patience 5. The Phase 2 recipe, fixed before any run, is AdamW~\cite{loshchilov2019decoupled} with learning rate
$2\times10^{-3}$ and weight decay $10^{-2}$ (not applied to biases and batch-normalisation parameters), one epoch
of linear warm-up followed by cosine decay to zero over 30 epochs~\cite{loshchilov2017sgdr}, batch size 128, label
smoothing 0.1, dropout 0.3 before the classifier, and the augmentation of Section~\ref{sec:p2data}. Early stopping
monitors the validation cross-entropy without label smoothing (patience 8) and the best weights are restored. On
the GPU the forward pass uses mixed precision. The seed fixes the initial weights, the batch order and the
augmentation draws.
